\documentclass[12pt,a4paper]{article}

%% ── Encoding & language ──────────────────────────────────────────────────────
\usepackage[T1]{fontenc}
\usepackage[utf8]{inputenc}
\usepackage[english]{babel}

%% ── Page geometry ────────────────────────────────────────────────────────────
\usepackage[top=2.5cm,bottom=2.5cm,left=3cm,right=3cm]{geometry}

%% ── Typography ───────────────────────────────────────────────────────────────
\usepackage[protrusion=false,expansion=false]{microtype}
\usepackage{parskip}
\setlength{\parindent}{0pt}

%% ── Math ─────────────────────────────────────────────────────────────────────
\usepackage{amsmath}
\usepackage{amssymb}

%% ── Tables ───────────────────────────────────────────────────────────────────
\usepackage{booktabs}
\usepackage{longtable}
\usepackage{array}
\newcolumntype{L}[1]{>{\raggedright\arraybackslash}p{#1}}

%% ── Colours ──────────────────────────────────────────────────────────────────
\usepackage{xcolor}
\definecolor{codebg}{RGB}{248,248,248}
\definecolor{kwcolor}{RGB}{0,0,180}
\definecolor{strcolor}{RGB}{163,21,21}
\definecolor{cmtcolor}{RGB}{0,128,0}
\definecolor{iricolor}{RGB}{0,100,160}
\definecolor{ruleblue}{RGB}{30,80,160}

%% ── Code listings ────────────────────────────────────────────────────────────
\usepackage{listings}

\lstdefinelanguage{turtle}{
  alsoletter={:,-,@},
  keywords={a, @prefix, @base, true, false},
  keywordstyle=\color{kwcolor}\bfseries,
  sensitive=true,
  comment=[l]{\#},
  commentstyle=\color{cmtcolor}\itshape,
  morestring=[b]",
  stringstyle=\color{strcolor},
  morestring=[s]{<}{>},
}

\lstdefinestyle{turtlestyle}{
  language=turtle,
  basicstyle=\ttfamily\scriptsize,
  backgroundcolor=\color{codebg},
  frame=single,
  framerule=0.4pt,
  rulecolor=\color{gray!50},
  numbers=none,
  columns=flexible,
  breaklines=true,
  breakatwhitespace=false,
  breakindent=12pt,
  postbreak=\mbox{\textcolor{gray}{$\hookrightarrow$}\space},
  showstringspaces=false,
  tabsize=2,
  xleftmargin=6pt,
  xrightmargin=6pt,
  aboveskip=6pt,
  belowskip=6pt,
}

\lstdefinestyle{pythonstyle}{
  language=Python,
  basicstyle=\ttfamily\scriptsize,
  backgroundcolor=\color{codebg},
  keywordstyle=\color{kwcolor}\bfseries,
  stringstyle=\color{strcolor},
  commentstyle=\color{cmtcolor}\itshape,
  frame=single,
  framerule=0.4pt,
  rulecolor=\color{gray!50},
  numbers=left,
  numberstyle=\tiny\color{gray},
  numbersep=8pt,
  columns=flexible,
  breaklines=true,
  breakatwhitespace=false,
  breakindent=12pt,
  postbreak=\mbox{\textcolor{gray}{$\hookrightarrow$}\space},
  showstringspaces=false,
  tabsize=4,
  xleftmargin=18pt,
  xrightmargin=6pt,
  aboveskip=6pt,
  belowskip=6pt,
}

\lstdefinestyle{sparqlstyle}{
  basicstyle=\ttfamily\scriptsize,
  backgroundcolor=\color{codebg},
  frame=single,
  framerule=0.4pt,
  rulecolor=\color{gray!50},
  numbers=none,
  columns=flexible,
  breaklines=true,
  breakatwhitespace=false,
  breakindent=12pt,
  postbreak=\mbox{\textcolor{gray}{$\hookrightarrow$}\space},
  showstringspaces=false,
  tabsize=2,
  xleftmargin=6pt,
  xrightmargin=6pt,
  aboveskip=6pt,
  belowskip=6pt,
}

%% ── Hyperlinks ───────────────────────────────────────────────────────────────
\usepackage{hyperref}
\hypersetup{
  colorlinks=true,
  linkcolor=ruleblue,
  urlcolor=ruleblue,
  citecolor=ruleblue,
  pdftitle={gUFO SHACL Validation -- Reference Documentation},
  pdfauthor={CDO-Shapes-gufo project},
  pdfsubject={gUFO SHACL validation shapes},
}

%% ── Headers and footers ──────────────────────────────────────────────────────
\usepackage{fancyhdr}
\pagestyle{fancy}
\fancyhf{}
\fancyhead[L]{\small\textit{gUFO SHACL Validation Reference}}
\fancyhead[R]{\small\textit{CDO-Shapes-gufo}}
\fancyfoot[C]{\small\thepage}
\renewcommand{\headrulewidth}{0.4pt}

%% ── Section formatting ───────────────────────────────────────────────────────
\usepackage{titlesec}
\titleformat{\section}{\large\bfseries\color{ruleblue}}{{\thesection}}{1em}{}[\titlerule]
\titleformat{\subsection}{\normalsize\bfseries\color{ruleblue}}{\thesubsection}{1em}{}
\titleformat{\subsubsection}{\normalsize\bfseries}{\thesubsubsection}{1em}{}

%% ── Custom commands ──────────────────────────────────────────────────────────
%% Using \nolinkurl (from hyperref) allows typewriter text to break at
%% special characters (:, ., -, _) preventing overfull hboxes in body text.
\newcommand{\ruleid}[1]{\textbf{\textsf{#1}}}
\newcommand{\shapename}[1]{\nolinkurl{#1}}
\newcommand{\gufoclass}[1]{\nolinkurl{gufo:#1}}
\newcommand{\gufoprop}[1]{\nolinkurl{gufo:#1}}
\newcommand{\shprop}[1]{\nolinkurl{sh:#1}}

%% ── Bibliography ─────────────────────────────────────────────────────────────
\usepackage{natbib}
\bibliographystyle{plainnat}

%% ── Overflow prevention ──────────────────────────────────────────────────────
%% Allow extra inter-word stretch as last resort before declaring overfull hbox.
\setlength{\emergencystretch}{3em}
%% Allow url package to break long URLs/paths at more characters.
\PassOptionsToPackage{hyphens}{url}

%% ─────────────────────────────────────────────────────────────────────────────
\begin{document}

\begin{titlepage}
  \centering
  \vspace*{3cm}
  {\Huge\bfseries\color{ruleblue} gUFO SHACL Validation\\[0.4em]
   Reference Documentation}\\[2em]
  {\large CDO-Shapes-gufo project}\\[1em]
  {\large\url{https://github.com/Cyber-Domain-Ontology/CDO-Shapes-gufo}}\\[3em]
  \hrule height 0.4pt
  \vspace{1em}
  {\normalsize
    Based on the gentle Unified Foundational Ontology (gUFO)\\
    Validation using SHACL via \texttt{pyshacl}
  }
  \vfill
  {\small\today}
\end{titlepage}

\tableofcontents
\newpage

%% ═══════════════════════════════════════════════════════════════════════════
\section{Purpose and Scope}
%% ═══════════════════════════════════════════════════════════════════════════

This document is the authoritative reference for the SHACL validation shapes
defined in the \texttt{CDO-Shapes-gufo} repository for ontologies based on the
\textbf{gentle Unified Foundational Ontology (gUFO)} \citep{gufo2022}.
It covers:

\begin{itemize}
  \item Which gUFO rules are enforced and how each rule is decomposed into
        independently verifiable constraints.
  \item The formal Description Logic (DL) specification of every constraint,
        making the mathematical meaning of each test unambiguous.
  \item The SHACL shapes, test exemplars, and \texttt{pytest} functions that
        realise each constraint.
\end{itemize}

The validation shapes are defined in \texttt{shapes/sh-gufo.ttl} and applied
using \texttt{pyshacl} \citep{pyshacl}. The test suite resides in the
\texttt{tests/} directory. This document covers the twelve gUFO rules whose
SHACL coverage required correction or extension; fully covered rules are
documented separately.

%% ═══════════════════════════════════════════════════════════════════════════
\section{Repository Structure}
%% ═══════════════════════════════════════════════════════════════════════════

\begin{lstlisting}[style=sparqlstyle]
CDO-Shapes-gufo/
  shapes/
    sh-gufo.ttl                  % SHACL shape definitions
  tests/
    exemplars.ttl                % Valid individuals exercising each shape
    exemplars_XFAIL.ttl          % Deliberately invalid individuals
    exemplars_XFAIL_validation.ttl  % Expected SHACL violation reports
    test_exemplar_coverage.py    % pytest test suite
    sh-qc.ttl                    % Quality-check shapes for test data
  dependencies/
    formatted-gufo.ttl           % gUFO ontology (OWL 2 DL)
\end{lstlisting}

%% ═══════════════════════════════════════════════════════════════════════════
\section{Validation Pipeline}
%% ═══════════════════════════════════════════════════════════════════════════

Validation proceeds in four sequential stages:

\begin{enumerate}
  \item \textbf{Entailment.} \texttt{src/entail.py} adds RDFS entailments to
    \texttt{exemplars.ttl} (subclass, subproperty, domain, range) so that
    \texttt{pyshacl} sees inherited types. This step applies only to
    \texttt{exemplars.ttl}, not to \texttt{exemplars\_XFAIL.ttl}. Failure test
    data must therefore declare all types explicitly.

  \item \textbf{Success validation.} \texttt{pyshacl} validates the expanded
    exemplars against \texttt{sh-gufo.ttl} and expects \nolinkurl{sh:conforms
    true}. Any violation indicates a regression in the shapes.

  \item \textbf{Failure validation.} \texttt{pyshacl} validates
    \texttt{exemplars\_XFAIL.ttl} and expects \nolinkurl{sh:conforms false}. The
    resulting violation report is saved as
    \texttt{exemplars\_XFAIL\_validation.ttl}.

  \item \textbf{pytest.} \texttt{test\_exemplar\_coverage.py} verifies two
    things:
    \begin{itemize}
      \item \texttt{test\_exemplar\_coverage()} confirms that every shape
        target class and property path has at least one exemplar individual
        in \texttt{exemplars.ttl}.
      \item Per-rule functions read \texttt{exemplars\_XFAIL\_validation.ttl}
        and assert that exactly the expected violations are present,
        identified by focus node, source shape, and constraint component.
    \end{itemize}
\end{enumerate}

%% ═══════════════════════════════════════════════════════════════════════════
\section{Literature and gUFO Rule Sources}
%% ═══════════════════════════════════════════════════════════════════════════

The rules documented here are derived from the following sources, listed in
order of authority:

\begin{longtable}{L{6cm} L{8cm}}
  \toprule
  \textbf{Reference} & \textbf{Scope} \\
  \midrule
  \endhead
  \texttt{formatted-gufo.ttl} --- the gUFO OWL~2 DL ontology
    & Authoritative axiomatisation \\[4pt]
  \citet{gufo2022} --- \textit{UFO: Unified Foundational Ontology},
    Applied Ontology 17(1)
    & Complete formal specification of all micro-theories \\[4pt]
  \citet{guizzardi2005} --- \textit{Ontological Foundations for Structural
    Conceptual Models}, PhD thesis, University of Twente
    & Original axiomatisation of UFO-A (endurants, aspects, relations) \\[4pt]
  \citet{guizzardi2018} --- \textit{Endurant Types in Ontology-Driven
    Conceptual Modeling: Towards OntoUML~2.0}, ER~2018
    & Endurant type taxonomy and constraints \\[4pt]
  \citet{fonseca2019} --- \textit{Relations in Ontology-Driven Conceptual
    Modeling}, ER~2019
    & Material, comparative, relator, and extrinsic mode constraints \\[4pt]
  \citet{carvalho2017} --- \textit{Multi-level ontology-based conceptual
    modeling}, DKE~109
    & \texttt{partitions} and \texttt{categorizes} powertype patterns \\[4pt]
  \citet{sales2015} --- \textit{Empirically Uncovered Anti-Patterns}, DKE~99
    & RelRig, FreeRole, MultDep anti-patterns \\[4pt]
  \citet{sales2019} --- \textit{Taxonomic Structures}, CEUR Vol-2519
    & MixIden, MixRig, GSRig anti-patterns \\
  \bottomrule
\end{longtable}

%% ═══════════════════════════════════════════════════════════════════════════
\section{Rule Documentation}
%% ═══════════════════════════════════════════════════════════════════════════

Each rule is documented in the following four-part structure:

\begin{description}
  \item[(i) Rule description] Natural language statement of the rule, its
    ontological rationale, and its source in the gUFO literature.
  \item[(ii) Test aspects] Decomposition into independently verifiable
    aspects, each with a unique identifier and description.
  \item[(iii) DL expressions] Formal Description Logic specification of each
    aspect, using standard DL notation: $\sqsubseteq$ for subsumption,
    $\sqcap$/$\sqcup$ for conjunction/disjunction, $\exists R.C$ and
    $\forall R.C$ for existential and universal restrictions, $\leq n\,R.C$
    and $\geq n\,R.C$ for qualified cardinality restrictions, and $\top$ for
    the universal concept. Inverse properties are written $R^-$.
  \item[(iv) Implementation] The SHACL shape(s), test exemplars (both valid
    and deliberately invalid), and \texttt{pytest} function(s).
\end{description}

\bigskip

%% ───────────────────────────────────────────────────────────────────────────
\subsection{Rule A1 --- Every Aspect Inheres in Exactly One ConcreteIndividual}
\label{sec:rule-a1}
%% ───────────────────────────────────────────────────────────────────────────

\subsubsection*{(i) Rule Description}

In gUFO a \gufoclass{Aspect} is an endurant that is existentially dependent on
another concrete individual---its \emph{bearer}---through the relation of
\emph{inherence}, expressed by \gufoprop{inheresIn}. The property is declared
functional, asymmetric, and irreflexive with domain \gufoclass{Aspect} and range
\gufoclass{ConcreteIndividual}. Both \gufoclass{IntrinsicAspect} and
\gufoclass{ExtrinsicMode} carry an \nolinkurl{owl:qualifiedCardinality "1"}
restriction on \texttt{inheresIn}, expressing that every aspect inheres in
\emph{exactly one} concrete individual---neither zero nor more than one.

\textbf{Sources:} \citet{guizzardi2005} ch.~4; \texttt{formatted-gufo.ttl}
\nolinkurl{owl:qualifiedCardinality "1"} on \texttt{IntrinsicAspect} and
\texttt{ExtrinsicMode}.

\subsubsection*{(ii) Test Aspects}

\begin{tabular}{lL{10cm}}
  \toprule
  \textbf{ID} & \textbf{Aspect} \\
  \midrule
  A1.1 & Upper bound: an \gufoclass{Aspect} may not inhere in more than one
         \gufoclass{ConcreteIndividual} \\
  A1.2 & Lower bound: every \gufoclass{Aspect} must inhere in at least one
         \gufoclass{ConcreteIndividual} \\
  A1.3 & Range: the object of \gufoprop{inheresIn} must be a
         \gufoclass{ConcreteIndividual} \\
  A1.4 & Domain: only an \gufoclass{Aspect} may be the subject of
         \gufoprop{inheresIn} \\
  \bottomrule
\end{tabular}

\subsubsection*{(iii) DL Expressions}

\textbf{A1.1 --- Upper bound:}
\begin{equation*}
  \mathsf{Aspect} \sqsubseteq {\leq}1\;\mathsf{inheresIn}.\mathsf{ConcreteIndividual}
\end{equation*}

\textbf{A1.2 --- Lower bound:}
\begin{equation*}
  \mathsf{Aspect} \sqsubseteq {\geq}1\;\mathsf{inheresIn}.\mathsf{ConcreteIndividual}
\end{equation*}

\textbf{Combined (full gUFO axiom):}
\begin{equation*}
  \mathsf{Aspect} \sqsubseteq {=}1\;\mathsf{inheresIn}.\mathsf{ConcreteIndividual}
\end{equation*}

\textbf{A1.3 --- Range:}
\begin{equation*}
  \top \sqsubseteq \forall\,\mathsf{inheresIn}.\mathsf{ConcreteIndividual}
\end{equation*}

\textbf{A1.4 --- Domain:}
\begin{equation*}
  \exists\,\mathsf{inheresIn}.\top \sqsubseteq \mathsf{Aspect}
\end{equation*}

\subsubsection*{(iv) Implementation}

\paragraph{SHACL shapes.}
\shapename{sh-gufo:Aspect-shape} enforces the lower bound (A1.2).
The upper bound (A1.1) is enforced by \shapename{sh-gufo:ExtrinsicMode-shape}
and \shapename{sh-gufo:IntrinsicAspect-shape} through
\shprop{qualifiedMaxCount~1}.

\begin{lstlisting}[style=turtlestyle,caption={sh-gufo:Aspect-shape with minCount~1}]
sh-gufo:Aspect-shape
    a sh:NodeShape ;
    rdfs:comment "rdfs:seeAlso links other shapes that review this class."@en ;
    rdfs:seeAlso sh-gufo:Aspect-disjointWith-Object-shape ;
    sh:property [
        a sh:PropertyShape ;
        sh:message "Every gufo:Aspect must inhere in exactly one ConcreteIndividual
            ( https://nemo-ufes.github.io/gufo/#inheresIn )."@en ;
        sh:minCount 1 ;
        sh:path gufo:inheresIn ;
    ] ;
    sh:targetClass gufo:Aspect ;
    sh:xone (
        [ a sh:NodeShape ; sh:class gufo:ExtrinsicAspect ; ]
        [ a sh:NodeShape ; sh:class gufo:IntrinsicAspect ; ]
    ) ;
    .

sh-gufo:inheresIn-objects-shape
    a sh:NodeShape ;
    sh:class gufo:ConcreteIndividual ;
    sh:nodeKind sh:BlankNodeOrIRI ;
    sh:targetObjectsOf gufo:inheresIn ;
    .

sh-gufo:inheresIn-subjects-shape
    a sh:NodeShape ;
    sh:class gufo:Aspect ;
    sh:property [
        a sh:PropertyShape ;
        sh:maxCount "1"^^xsd:integer ;
        sh:path gufo:inheresIn ;
    ] ;
    sh:targetSubjectsOf gufo:inheresIn ;
    .
\end{lstlisting}

\paragraph{Success test (\texttt{tests/exemplars.ttl}).}

\begin{lstlisting}[style=turtlestyle]
kb:Aspect-c6c0dc61-2ced-498d-821c-443761c87de0
    a gufo:ExtrinsicMode ;
    gufo:inheresIn kb:ConcreteIndividual-1c404c3c-ecf3-4946-8130-e40548bef69f ;
    rdfs:seeAlso
        sh-gufo:ExtrinsicMode-shape ,
        sh-gufo:inheresIn-subjects-shape ;
    .

kb:ConcreteIndividual-1c404c3c-ecf3-4946-8130-e40548bef69f
    a gufo:Object ;
    rdfs:seeAlso sh-gufo:inheresIn-objects-shape ;
    .
\end{lstlisting}

\paragraph{Failure test (\texttt{tests/exemplars\_XFAIL.ttl}) --- A1.2.}

\begin{lstlisting}[style=turtlestyle]
kb:ExtrinsicMode-8f4c2a1d-3b7e-4f9a-bc12-5d6e7f8a9b0c
    a gufo:ExtrinsicMode ;
    rdfs:comment "This node is expected to trigger a validation error for
        being a gufo:Aspect without any gufo:inheresIn triple."@en ;
    rdfs:seeAlso sh-gufo:Aspect-shape ;
    .
\end{lstlisting}

\paragraph{pytest (\texttt{tests/test\_exemplar\_coverage.py}).}

\begin{lstlisting}[style=pythonstyle]
def test_exemplar_xfail_validation_aspect_inheres_in() -> None:
    """
    Verifies A1.2: every gufo:Aspect must inhere in at least one
    ConcreteIndividual (sh:minCount 1 on gufo:inheresIn in Aspect-shape).
    An ExtrinsicMode without any inheresIn triple must trigger a
    MinCountConstraintComponent violation on Aspect-shape.
    """
    validation_graph = Graph()
    validation_graph.parse("exemplars_XFAIL_validation.ttl")
    ns_kb = Namespace("http://example.org/kb/")
    n_focus_nodes: Set[URIRef] = set()
    for result in validation_graph.query("""\
PREFIX sh: <http://www.w3.org/ns/shacl#>
PREFIX sh-gufo: <http://example.org/shapes/sh-gufo/>
SELECT ?nAspect
WHERE {
  ?nValidationResult
    a sh:ValidationResult ;
    sh:sourceShape sh-gufo:Aspect-shape ;
    sh:sourceConstraintComponent sh:MinCountConstraintComponent ;
    sh:focusNode ?nAspect ;
    .
}
"""):
        assert isinstance(result, ResultRow)
        assert isinstance(result[0], URIRef)
        n_focus_nodes.add(result[0])
    assert n_focus_nodes == {
        ns_kb["ExtrinsicMode-8f4c2a1d-3b7e-4f9a-bc12-5d6e7f8a9b0c"],
    }
\end{lstlisting}

%% ───────────────────────────────────────────────────────────────────────────
\subsection{Rule A2 --- Every ExtrinsicMode Must Externally Depend on At Least One Endurant}
\label{sec:rule-a2}
%% ───────────────────────────────────────────────────────────────────────────

\subsubsection*{(i) Rule Description}

A \gufoclass{ExtrinsicMode} is an extrinsic aspect that inheres in one concrete
individual (its bearer) but also \emph{externally depends} on at least one other
endurant mereologically disjoint from the bearer, expressed via
\gufoprop{externallyDependsOn}. The gUFO OWL implementation restricts
\gufoclass{ExtrinsicMode} with \nolinkurl{owl:someValuesFrom gufo:ConcreteIndividual}
on \texttt{externallyDependsOn}. The property is also declared
\nolinkurl{owl:IrreflexiveProperty}.

\textbf{Sources:} \citet{guizzardi2005}; \texttt{formatted-gufo.ttl};
gUFO documentation \S2.7.

\subsubsection*{(ii) Test Aspects}

\begin{tabular}{lL{10cm}}
  \toprule
  \textbf{ID} & \textbf{Aspect} \\
  \midrule
  A2.1 & Existential obligation: every \gufoclass{ExtrinsicMode} must have at
         least one \gufoprop{externallyDependsOn} triple \\
  A2.2 & Domain: only an \gufoclass{ExtrinsicMode} may be the subject of
         \gufoprop{externallyDependsOn} \\
  A2.3 & Range: the object of \gufoprop{externallyDependsOn} must be an
         \gufoclass{Endurant} \\
  \bottomrule
\end{tabular}

\subsubsection*{(iii) DL Expressions}

\textbf{A2.1 --- Existential obligation:}
\begin{equation*}
  \mathsf{ExtrinsicMode} \sqsubseteq {\geq}1\;\mathsf{externallyDependsOn}.\mathsf{Endurant}
\end{equation*}

\textbf{A2.2 --- Domain:}
\begin{equation*}
  \exists\,\mathsf{externallyDependsOn}.\top \sqsubseteq \mathsf{ExtrinsicMode}
\end{equation*}

\textbf{A2.3 --- Range:}
\begin{equation*}
  \top \sqsubseteq \forall\,\mathsf{externallyDependsOn}.\mathsf{Endurant}
\end{equation*}

\subsubsection*{(iv) Implementation}

\paragraph{SHACL shapes.}

\begin{lstlisting}[style=turtlestyle]
sh-gufo:ExtrinsicMode-shape
    a sh:NodeShape ;
    sh:property
        [
            a sh:PropertyShape ;
            sh:message "Every gufo:ExtrinsicMode must externally depend on at
                least one Endurant ( https://nemo-ufes.github.io/gufo/#externallyDependsOn )."@en ;
            sh:minCount 1 ;
            sh:path gufo:externallyDependsOn ;
        ] ,
        [
            a sh:PropertyShape ;
            sh:path gufo:inheresIn ;
            sh:qualifiedMaxCount "1"^^xsd:integer ;
            sh:qualifiedValueShape [
                a sh:NodeShape ;
                sh:class gufo:ConcreteIndividual ;
            ] ;
        ] ;
    sh:targetClass gufo:ExtrinsicMode ;
    .

sh-gufo:externallyDependsOn-subjects-shape
    a sh:NodeShape ;
    sh:class gufo:ExtrinsicMode ;
    sh:targetSubjectsOf gufo:externallyDependsOn ;
    .

sh-gufo:externallyDependsOn-objects-shape
    a sh:NodeShape ;
    sh:class gufo:Endurant ;
    sh:nodeKind sh:BlankNodeOrIRI ;
    sh:targetObjectsOf gufo:externallyDependsOn ;
    .
\end{lstlisting}

\paragraph{Success test.}

\begin{lstlisting}[style=turtlestyle]
kb:ExtrinsicMode-01964f52-6337-4703-80d2-5903d7635901
    a gufo:ExtrinsicMode ;
    gufo:externallyDependsOn kb:Endurant-2617295a-5922-4d06-90ce-83845c2914af ;
    rdfs:seeAlso sh-gufo:externallyDependsOn-subjects-shape ;
    .

kb:Endurant-2617295a-5922-4d06-90ce-83845c2914af
    a gufo:Object ;
    rdfs:seeAlso sh-gufo:externallyDependsOn-objects-shape ;
    .
\end{lstlisting}

\paragraph{Failure test --- A2.1.}

\begin{lstlisting}[style=turtlestyle]
kb:ExtrinsicMode-3c7d1e2f-5a6b-4c8d-9e0f-1a2b3c4d5e6f
    a gufo:ExtrinsicMode ;
    gufo:inheresIn kb:ConcreteIndividual-1c404c3c-ecf3-4946-8130-e40548bef69f ;
    rdfs:comment "This node is expected to trigger a validation error for
        being a gufo:ExtrinsicMode without any gufo:externallyDependsOn triple."@en ;
    rdfs:seeAlso sh-gufo:ExtrinsicMode-shape ;
    .
\end{lstlisting}

\paragraph{pytest.}

\begin{lstlisting}[style=pythonstyle]
def test_exemplar_xfail_validation_extrinsicmode_externally_depends_on() -> None:
    """
    Verifies A2.1: every gufo:ExtrinsicMode must externally depend on at
    least one Endurant (sh:minCount 1 on gufo:externallyDependsOn).
    """
    validation_graph = Graph()
    validation_graph.parse("exemplars_XFAIL_validation.ttl")
    ns_kb = Namespace("http://example.org/kb/")
    n_focus_nodes: Set[URIRef] = set()
    for result in validation_graph.query("""\
PREFIX sh: <http://www.w3.org/ns/shacl#>
PREFIX sh-gufo: <http://example.org/shapes/sh-gufo/>
SELECT ?nExtrinsicMode
WHERE {
  ?nValidationResult
    a sh:ValidationResult ;
    sh:sourceShape sh-gufo:ExtrinsicMode-shape ;
    sh:sourceConstraintComponent sh:MinCountConstraintComponent ;
    sh:focusNode ?nExtrinsicMode ;
    .
}
"""):
        assert isinstance(result, ResultRow)
        assert isinstance(result[0], URIRef)
        n_focus_nodes.add(result[0])
    assert n_focus_nodes == {
        ns_kb["ExtrinsicMode-3c7d1e2f-5a6b-4c8d-9e0f-1a2b3c4d5e6f"],
    }
\end{lstlisting}

%% ───────────────────────────────────────────────────────────────────────────
\subsection{Rule A3 --- TemporaryConstitutionSituation Has Exactly One concernsConstitutedEndurant}
\label{sec:rule-a3}
%% ───────────────────────────────────────────────────────────────────────────

\subsubsection*{(i) Rule Description}

A \gufoclass{TemporaryConstitutionSituation} represents a situation in which one
endurant temporarily constitutes another. It must be associated with exactly one
constituted endurant via \gufoprop{concernsConstitutedEndurant}, and must be
referenced by exactly one constituting endurant via the inverse of
\gufoprop{standsInQualifiedConstitution}. The gUFO OWL implementation declares
\nolinkurl{owl:qualifiedCardinality "1"} on both paths.

\textbf{Sources:} \texttt{formatted-gufo.ttl}; gUFO documentation \S2.1.

\subsubsection*{(ii) Test Aspects}

\begin{tabular}{lL{9.5cm}}
  \toprule
  \textbf{ID} & \textbf{Aspect} \\
  \midrule
  A3.1 & Upper bound: at most one \gufoprop{concernsConstitutedEndurant} \\
  A3.2 & Lower bound: at least one \gufoprop{concernsConstitutedEndurant} \\
  A3.3 & Range: object must be an \gufoclass{Endurant} \\
  A3.4 & Domain: subject must be a \gufoclass{TemporaryConstitutionSituation} \\
  A3.5 & Upper bound on inverse: at most one \gufoclass{Endurant} via
         \gufoprop{standsInQualifiedConstitution} \\
  \bottomrule
\end{tabular}

\subsubsection*{(iii) DL Expressions}

\textbf{A3.1 --- Upper bound:}
\begin{equation*}
  \mathsf{TCS} \sqsubseteq {\leq}1\;\mathsf{concernsConstitutedEndurant}.\mathsf{Object}
\end{equation*}

\textbf{A3.2 --- Lower bound:}
\begin{equation*}
  \mathsf{TCS} \sqsubseteq {\geq}1\;\mathsf{concernsConstitutedEndurant}.\top
\end{equation*}

\textbf{Combined:}
\begin{equation*}
  \mathsf{TCS} \sqsubseteq {=}1\;\mathsf{concernsConstitutedEndurant}.\mathsf{Object}
\end{equation*}

where $\mathsf{TCS}$ abbreviates $\mathsf{TemporaryConstitutionSituation}$.

\textbf{A3.3 --- Range:}
\begin{equation*}
  \top \sqsubseteq \forall\,\mathsf{concernsConstitutedEndurant}.\mathsf{Endurant}
\end{equation*}

\textbf{A3.4 --- Domain:}
\begin{equation*}
  \exists\,\mathsf{concernsConstitutedEndurant}.\top \sqsubseteq \mathsf{TCS}
\end{equation*}

\textbf{A3.5 --- Upper bound on inverse:}
\begin{equation*}
  \mathsf{TCS} \sqsubseteq {\leq}1\;\mathsf{standsInQualifiedConstitution}^{-}.\mathsf{Object}
\end{equation*}

\subsubsection*{(iv) Implementation}

\paragraph{SHACL shapes.}

\begin{lstlisting}[style=turtlestyle]
sh-gufo:TemporaryConstitutionSituation-shape
    a sh:NodeShape ;
    sh:property
        [
            a sh:PropertyShape ;
            sh:message "Every gufo:TemporaryConstitutionSituation must concern
                exactly one constituted Endurant."@en ;
            sh:minCount 1 ;
            sh:path gufo:concernsConstitutedEndurant ;
        ] ,
        [
            a sh:PropertyShape ;
            sh:path gufo:concernsConstitutedEndurant ;
            sh:qualifiedMaxCount "1"^^xsd:integer ;
            sh:qualifiedValueShape [ a sh:NodeShape ; sh:class gufo:Object ; ] ;
        ] ,
        [
            a sh:PropertyShape ;
            sh:path [ sh:inversePath gufo:standsInQualifiedConstitution ] ;
            sh:qualifiedMaxCount "1"^^xsd:integer ;
            sh:qualifiedValueShape [ a sh:NodeShape ; sh:class gufo:Object ; ] ;
        ] ;
    sh:targetClass gufo:TemporaryConstitutionSituation ;
    .
\end{lstlisting}

\paragraph{Success tests.}

\begin{lstlisting}[style=turtlestyle]
kb:TemporaryConstitutionSituation-490fa048-fc2a-4434-8376-80579af09d88
    a gufo:TemporaryConstitutionSituation ;
    gufo:concernsConstitutedEndurant kb:Endurant-11b7d0cf-3a5e-40da-a79c-a6b9c243ee28 ;
    rdfs:seeAlso
        sh-gufo:TemporaryConstitutionSituation-shape ,
        sh-gufo:concernsConstitutedEndurant-subjects-shape ;
    .

kb:TemporaryConstitutionSituation-932e7a34-1d58-4e8b-8265-9f8f62bb0c34
    a gufo:TemporaryConstitutionSituation ;
    rdfs:seeAlso sh-gufo:standsInQualifiedConstitution-objects-shape ;
    .
\end{lstlisting}

\paragraph{Failure test --- A3.2.}

\begin{lstlisting}[style=turtlestyle]
kb:TemporaryConstitutionSituation-4d8e2f3a-6b7c-4d9e-0f1a-2b3c4d5e6f7a
    a gufo:TemporaryConstitutionSituation ;
    rdfs:comment "This node is expected to trigger a validation error for being
        a gufo:TemporaryConstitutionSituation without any
        gufo:concernsConstitutedEndurant triple."@en ;
    rdfs:seeAlso sh-gufo:TemporaryConstitutionSituation-shape ;
    .
\end{lstlisting}

\paragraph{pytest.}

\begin{lstlisting}[style=pythonstyle]
def test_exemplar_xfail_validation_temporary_constitution_situation() -> None:
    """Verifies A3.2: minCount 1 on concernsConstitutedEndurant."""
    validation_graph = Graph()
    validation_graph.parse("exemplars_XFAIL_validation.ttl")
    ns_kb = Namespace("http://example.org/kb/")
    n_focus_nodes: Set[URIRef] = set()
    for result in validation_graph.query("""\
PREFIX sh: <http://www.w3.org/ns/shacl#>
PREFIX sh-gufo: <http://example.org/shapes/sh-gufo/>
SELECT ?nSituation
WHERE {
  ?nValidationResult
    a sh:ValidationResult ;
    sh:sourceShape sh-gufo:TemporaryConstitutionSituation-shape ;
    sh:sourceConstraintComponent sh:MinCountConstraintComponent ;
    sh:focusNode ?nSituation ;
    .
}
"""):
        assert isinstance(result, ResultRow)
        assert isinstance(result[0], URIRef)
        n_focus_nodes.add(result[0])
    assert n_focus_nodes == {
        ns_kb["TemporaryConstitutionSituation-4d8e2f3a-6b7c-4d9e-0f1a-2b3c4d5e6f7a"],
    }
\end{lstlisting}

%% ───────────────────────────────────────────────────────────────────────────
\subsection{Rule A4 --- TemporaryInstantiationSituation Has Exactly One concernsNonRigidType}
\label{sec:rule-a4}
%% ───────────────────────────────────────────────────────────────────────────

\subsubsection*{(i) Rule Description}

A \gufoclass{TemporaryInstantiationSituation} captures a situation in which an
endurant temporarily instantiates a \gufoclass{NonRigidType}. It must be
associated with exactly one \gufoclass{NonRigidType} via
\gufoprop{concernsNonRigidType}, and must be referenced by exactly one endurant
via the inverse of \gufoprop{standsInQualifiedInstantiation}. The gUFO OWL
implementation declares \nolinkurl{owl:qualifiedCardinality "1"} on both paths.

\textbf{Sources:} \texttt{formatted-gufo.ttl}; gUFO documentation \S2.9.

\subsubsection*{(ii) Test Aspects}

\begin{tabular}{lL{9.5cm}}
  \toprule
  \textbf{ID} & \textbf{Aspect} \\
  \midrule
  A4.1 & Upper bound: at most one \gufoprop{concernsNonRigidType} \\
  A4.2 & Lower bound: at least one \gufoprop{concernsNonRigidType} \\
  A4.3 & Range: object must be a \gufoclass{NonRigidType} \\
  A4.4 & Domain: subject must be a \gufoclass{TemporaryInstantiationSituation} \\
  A4.5 & Upper bound on inverse: at most one \gufoclass{Endurant} via
         \gufoprop{standsInQualifiedInstantiation} \\
  \bottomrule
\end{tabular}

\subsubsection*{(iii) DL Expressions}

Let $\mathsf{TIS}$ abbreviate $\mathsf{TemporaryInstantiationSituation}$.

\textbf{A4.1:}
$\mathsf{TIS} \sqsubseteq {\leq}1\;\mathsf{concernsNonRigidType}.\mathsf{NonRigidType}$

\textbf{A4.2:}
$\mathsf{TIS} \sqsubseteq {\geq}1\;\mathsf{concernsNonRigidType}.\top$

\textbf{Combined:}
$\mathsf{TIS} \sqsubseteq {=}1\;\mathsf{concernsNonRigidType}.\mathsf{NonRigidType}$

\textbf{A4.3:}
$\top \sqsubseteq \forall\,\mathsf{concernsNonRigidType}.\mathsf{NonRigidType}$

\textbf{A4.4:}
$\exists\,\mathsf{concernsNonRigidType}.\top \sqsubseteq \mathsf{TIS}$

\textbf{A4.5:}
$\mathsf{TIS} \sqsubseteq {\leq}1\;\mathsf{standsInQualifiedInstantiation}^{-}.\mathsf{Endurant}$

\subsubsection*{(iv) Implementation}

\paragraph{SHACL shapes.}

\begin{lstlisting}[style=turtlestyle]
sh-gufo:TemporaryInstantiationSituation-shape
    a sh:NodeShape ;
    sh:property
        [
            a sh:PropertyShape ;
            sh:message "Every gufo:TemporaryInstantiationSituation must concern
                at least one NonRigidType."@en ;
            sh:minCount 1 ;
            sh:path gufo:concernsNonRigidType ;
        ] ,
        [
            a sh:PropertyShape ;
            sh:path gufo:concernsNonRigidType ;
            sh:qualifiedMaxCount "1"^^xsd:integer ;
            sh:qualifiedValueShape [ a sh:NodeShape ; sh:class gufo:NonRigidType ; ] ;
        ] ,
        [
            a sh:PropertyShape ;
            sh:path [ sh:inversePath gufo:standsInQualifiedInstantiation ] ;
            sh:qualifiedMaxCount "1"^^xsd:integer ;
            sh:qualifiedValueShape [ a sh:NodeShape ; sh:class gufo:Endurant ; ] ;
        ] ;
    sh:targetClass gufo:TemporaryInstantiationSituation ;
    .
\end{lstlisting}

\paragraph{Success tests.}

\begin{lstlisting}[style=turtlestyle]
kb:TemporaryInstantiationSituation-41702321-9550-458b-9438-f48d28d7e0c8
    a gufo:TemporaryInstantiationSituation ;
    gufo:concernsNonRigidType kb:NonRigidType-0a7a7dc6-842e-4542-a360-56c4f2f13308 ;
    rdfs:seeAlso
        sh-gufo:TemporaryInstantiationSituation-shape ,
        sh-gufo:concernsNonRigidType-subjects-shape ;
    .

kb:TemporaryInstantiationSituation-d50f30d2-d042-4003-9b81-5878e348bbc6
    a gufo:TemporaryInstantiationSituation ;
    rdfs:seeAlso sh-gufo:standsInQualifiedInstantiation-objects-shape ;
    .
\end{lstlisting}

\paragraph{Failure test --- A4.2.}

\begin{lstlisting}[style=turtlestyle]
kb:TemporaryInstantiationSituation-5e9f3a4b-7c8d-4e0f-1a2b-3c4d5e6f7a8b
    a gufo:TemporaryInstantiationSituation ;
    rdfs:comment "This node is expected to trigger a validation error for
        being a gufo:TemporaryInstantiationSituation without any
        gufo:concernsNonRigidType triple."@en ;
    rdfs:seeAlso sh-gufo:TemporaryInstantiationSituation-shape ;
    .
\end{lstlisting}

\paragraph{pytest.}

\begin{lstlisting}[style=pythonstyle]
def test_exemplar_xfail_validation_temporary_instantiation_situation() -> None:
    """Verifies A4.2: minCount 1 on concernsNonRigidType."""
    validation_graph = Graph()
    validation_graph.parse("exemplars_XFAIL_validation.ttl")
    ns_kb = Namespace("http://example.org/kb/")
    n_focus_nodes: Set[URIRef] = set()
    for result in validation_graph.query("""\
PREFIX sh: <http://www.w3.org/ns/shacl#>
PREFIX sh-gufo: <http://example.org/shapes/sh-gufo/>
SELECT ?nSituation
WHERE {
  ?nValidationResult
    a sh:ValidationResult ;
    sh:sourceShape sh-gufo:TemporaryInstantiationSituation-shape ;
    sh:sourceConstraintComponent sh:MinCountConstraintComponent ;
    sh:focusNode ?nSituation ;
    .
}
"""):
        assert isinstance(result, ResultRow)
        assert isinstance(result[0], URIRef)
        n_focus_nodes.add(result[0])
    assert n_focus_nodes == {
        ns_kb["TemporaryInstantiationSituation-5e9f3a4b-7c8d-4e0f-1a2b-3c4d5e6f7a8b"],
    }
\end{lstlisting}

%% ───────────────────────────────────────────────────────────────────────────
\subsection{Rule A5 --- TemporaryParthoodSituation Has Exactly One concernsTemporaryWhole}
\label{sec:rule-a5}
%% ───────────────────────────────────────────────────────────────────────────

\subsubsection*{(i) Rule Description}

A \gufoclass{TemporaryParthoodSituation} captures a situation in which an
endurant is temporarily part of another. It must be associated with exactly one
whole via \gufoprop{concernsTemporaryWhole}, and must be referenced by exactly
one endurant via the inverse of \gufoprop{standsInQualifiedParthood}. The gUFO
OWL implementation declares \nolinkurl{owl:qualifiedCardinality "1"} on both paths.

\textbf{Sources:} \texttt{formatted-gufo.ttl}; gUFO documentation \S2.2.

\subsubsection*{(ii) Test Aspects}

\begin{tabular}{lL{9.5cm}}
  \toprule
  \textbf{ID} & \textbf{Aspect} \\
  \midrule
  A5.1 & Upper bound: at most one \gufoprop{concernsTemporaryWhole} \\
  A5.2 & Lower bound: at least one \gufoprop{concernsTemporaryWhole} \\
  A5.3 & Range: object must be an \gufoclass{Endurant} \\
  A5.4 & Domain: subject must be a \gufoclass{TemporaryParthoodSituation} \\
  A5.5 & Upper bound on inverse: at most one \gufoclass{Endurant} via
         \gufoprop{standsInQualifiedParthood} \\
  \bottomrule
\end{tabular}

\subsubsection*{(iii) DL Expressions}

Let $\mathsf{TPS}$ abbreviate $\mathsf{TemporaryParthoodSituation}$.

\textbf{A5.1:}
$\mathsf{TPS} \sqsubseteq {\leq}1\;\mathsf{concernsTemporaryWhole}.\mathsf{Endurant}$

\textbf{A5.2:}
$\mathsf{TPS} \sqsubseteq {\geq}1\;\mathsf{concernsTemporaryWhole}.\top$

\textbf{Combined:}
$\mathsf{TPS} \sqsubseteq {=}1\;\mathsf{concernsTemporaryWhole}.\mathsf{Endurant}$

\textbf{A5.3:}
$\top \sqsubseteq \forall\,\mathsf{concernsTemporaryWhole}.\mathsf{Endurant}$

\textbf{A5.4:}
$\exists\,\mathsf{concernsTemporaryWhole}.\top \sqsubseteq \mathsf{TPS}$

\textbf{A5.5:}
$\mathsf{TPS} \sqsubseteq {\leq}1\;\mathsf{standsInQualifiedParthood}^{-}.\mathsf{Endurant}$

\subsubsection*{(iv) Implementation}

\paragraph{SHACL shapes.}

\begin{lstlisting}[style=turtlestyle]
sh-gufo:TemporaryParthoodSituation-shape
    a sh:NodeShape ;
    sh:property
        [
            a sh:PropertyShape ;
            sh:message "Every gufo:TemporaryParthoodSituation must concern at
                least one temporary whole."@en ;
            sh:minCount 1 ;
            sh:path gufo:concernsTemporaryWhole ;
        ] ,
        [
            a sh:PropertyShape ;
            sh:path gufo:concernsTemporaryWhole ;
            sh:qualifiedMaxCount "1"^^xsd:integer ;
            sh:qualifiedValueShape [ a sh:NodeShape ; sh:class gufo:Endurant ; ] ;
        ] ,
        [
            a sh:PropertyShape ;
            sh:path [ sh:inversePath gufo:standsInQualifiedParthood ] ;
            sh:qualifiedMaxCount "1"^^xsd:integer ;
            sh:qualifiedValueShape [ a sh:NodeShape ; sh:class gufo:Endurant ; ] ;
        ] ;
    sh:targetClass gufo:TemporaryParthoodSituation ;
    .
\end{lstlisting}

\paragraph{Success tests.}

\begin{lstlisting}[style=turtlestyle]
kb:TemporaryParthoodSituation-4292a36e-659c-4644-a197-a0e84dfb2394
    a gufo:TemporaryParthoodSituation ;
    gufo:concernsTemporaryWhole kb:Endurant-4053ed7d-05c9-44ce-869f-2745d84a40ce ;
    rdfs:seeAlso
        sh-gufo:TemporaryParthoodSituation-shape ,
        sh-gufo:concernsTemporaryWhole-subjects-shape ;
    .

kb:TemporaryParthoodSituation-bdfd62c0-4188-4cc2-a823-2664b8478fe1
    a gufo:TemporaryParthoodSituation ;
    rdfs:seeAlso sh-gufo:standsInQualifiedParthood-objects-shape ;
    .
\end{lstlisting}

\paragraph{Failure test --- A5.2.}

\begin{lstlisting}[style=turtlestyle]
kb:TemporaryParthoodSituation-6f0a4b5c-8d9e-4f1a-2b3c-4d5e6f7a8b9c
    a gufo:TemporaryParthoodSituation ;
    rdfs:comment "This node is expected to trigger a validation error for
        being a gufo:TemporaryParthoodSituation without any
        gufo:concernsTemporaryWhole triple."@en ;
    rdfs:seeAlso sh-gufo:TemporaryParthoodSituation-shape ;
    .
\end{lstlisting}

\paragraph{pytest.}

\begin{lstlisting}[style=pythonstyle]
def test_exemplar_xfail_validation_temporary_parthood_situation() -> None:
    """Verifies A5.2: minCount 1 on concernsTemporaryWhole."""
    validation_graph = Graph()
    validation_graph.parse("exemplars_XFAIL_validation.ttl")
    ns_kb = Namespace("http://example.org/kb/")
    n_focus_nodes: Set[URIRef] = set()
    for result in validation_graph.query("""\
PREFIX sh: <http://www.w3.org/ns/shacl#>
PREFIX sh-gufo: <http://example.org/shapes/sh-gufo/>
SELECT ?nSituation
WHERE {
  ?nValidationResult
    a sh:ValidationResult ;
    sh:sourceShape sh-gufo:TemporaryParthoodSituation-shape ;
    sh:sourceConstraintComponent sh:MinCountConstraintComponent ;
    sh:focusNode ?nSituation ;
    .
}
"""):
        assert isinstance(result, ResultRow)
        assert isinstance(result[0], URIRef)
        n_focus_nodes.add(result[0])
    assert n_focus_nodes == {
        ns_kb["TemporaryParthoodSituation-6f0a4b5c-8d9e-4f1a-2b3c-4d5e6f7a8b9c"],
    }
\end{lstlisting}

%% ───────────────────────────────────────────────────────────────────────────
\subsection{Rule A6 --- TemporaryRelationshipSituation Has One RelationshipType and At Least One Related Endurant}
\label{sec:rule-a6}
%% ───────────────────────────────────────────────────────────────────────────

\subsubsection*{(i) Rule Description}

A \gufoclass{TemporaryRelationshipSituation} reifies a temporary relationship
between endurants. It must be associated with exactly one
\gufoclass{RelationshipType} via \gufoprop{concernsRelationshipType}
(\nolinkurl{owl:qualifiedCardinality "1"}) and with at least one related endurant
via \gufoprop{concernsRelatedEndurant} (\nolinkurl{owl:someValuesFrom
gufo:Endurant}). It must also be referenced by exactly one endurant via the
inverse of \gufoprop{standsInQualifiedRelationship}.

\textbf{Sources:} \texttt{formatted-gufo.ttl}; gUFO documentation \S2.9.

\subsubsection*{(ii) Test Aspects}

\begin{tabular}{lL{9cm}}
  \toprule
  \textbf{ID} & \textbf{Aspect} \\
  \midrule
  A6.1 & Upper bound on \gufoprop{concernsRelationshipType}: at most one \\
  A6.2 & Lower bound on \gufoprop{concernsRelationshipType}: at least one \\
  A6.3 & Lower bound on \gufoprop{concernsRelatedEndurant}: at least one \\
  A6.4 & Upper bound on \gufoprop{concernsRelatedEndurant}: at most one \\
  A6.5 & Range of \gufoprop{concernsRelationshipType}: must be a \gufoclass{RelationshipType} \\
  A6.6 & Range of \gufoprop{concernsRelatedEndurant}: must be an \gufoclass{Endurant} \\
  A6.7 & Domain of \gufoprop{concernsRelationshipType}: must be a
         \gufoclass{TemporaryRelationshipSituation} \\
  A6.8 & Domain of \gufoprop{concernsRelatedEndurant}: must be a
         \gufoclass{TemporaryRelationshipSituation} \\
  A6.9 & Upper bound on inverse: at most one \gufoclass{Endurant} via
         \gufoprop{standsInQualifiedRelationship} \\
  \bottomrule
\end{tabular}

\subsubsection*{(iii) DL Expressions}

Let $\mathsf{TRS}$ abbreviate $\mathsf{TemporaryRelationshipSituation}$.

\begin{align*}
  \text{A6.1:}&\quad \mathsf{TRS} \sqsubseteq {\leq}1\;\mathsf{concernsRelationshipType}.\mathsf{RelationshipType} \\
  \text{A6.2:}&\quad \mathsf{TRS} \sqsubseteq {\geq}1\;\mathsf{concernsRelationshipType}.\top \\
  \text{A6.3:}&\quad \mathsf{TRS} \sqsubseteq {\geq}1\;\mathsf{concernsRelatedEndurant}.\mathsf{Endurant} \\
  \text{A6.4:}&\quad \mathsf{TRS} \sqsubseteq {\leq}1\;\mathsf{concernsRelatedEndurant}.\mathsf{Endurant} \\
  \text{A6.5:}&\quad \top \sqsubseteq \forall\,\mathsf{concernsRelationshipType}.\mathsf{RelationshipType} \\
  \text{A6.6:}&\quad \top \sqsubseteq \forall\,\mathsf{concernsRelatedEndurant}.\mathsf{Endurant} \\
  \text{A6.7:}&\quad \exists\,\mathsf{concernsRelationshipType}.\top \sqsubseteq \mathsf{TRS} \\
  \text{A6.8:}&\quad \exists\,\mathsf{concernsRelatedEndurant}.\top \sqsubseteq \mathsf{TRS} \\
  \text{A6.9:}&\quad \mathsf{TRS} \sqsubseteq {\leq}1\;\mathsf{standsInQualifiedRelationship}^{-}.\mathsf{Endurant}
\end{align*}

\subsubsection*{(iv) Implementation}

\paragraph{SHACL shapes.}

\begin{lstlisting}[style=turtlestyle]
sh-gufo:TemporaryRelationshipSituation-shape
    a sh:NodeShape ;
    sh:property
        [
            a sh:PropertyShape ;
            sh:message "Every gufo:TemporaryRelationshipSituation must concern
                at least one RelationshipType."@en ;
            sh:minCount 1 ;
            sh:path gufo:concernsRelationshipType ;
        ] ,
        [
            a sh:PropertyShape ;
            sh:path gufo:concernsRelationshipType ;
            sh:qualifiedMaxCount "1"^^xsd:integer ;
            sh:qualifiedValueShape [ a sh:NodeShape ; sh:class gufo:RelationshipType ; ] ;
        ] ,
        [
            a sh:PropertyShape ;
            sh:message "Every gufo:TemporaryRelationshipSituation must concern
                at least one related Endurant."@en ;
            sh:minCount 1 ;
            sh:path gufo:concernsRelatedEndurant ;
        ] ,
        [
            a sh:PropertyShape ;
            sh:path gufo:concernsRelatedEndurant ;
            sh:qualifiedMaxCount "1"^^xsd:integer ;
            sh:qualifiedValueShape [ a sh:NodeShape ; sh:class gufo:Endurant ; ] ;
        ] ,
        [
            a sh:PropertyShape ;
            sh:path [ sh:inversePath gufo:standsInQualifiedRelationship ] ;
            sh:qualifiedMaxCount "1"^^xsd:integer ;
            sh:qualifiedValueShape [ a sh:NodeShape ; sh:class gufo:Endurant ; ] ;
        ] ;
    sh:targetClass gufo:TemporaryRelationshipSituation ;
    .
\end{lstlisting}

\paragraph{Success tests.}

\begin{lstlisting}[style=turtlestyle]
kb:TemporaryRelationshipSituation-f874b7a3-5add-49f4-895c-9fbe42b40095
    a gufo:TemporaryRelationshipSituation ;
    gufo:concernsRelationshipType kb:RelationshipType-0813ed87-2f47-4a1b-a063-834b467c3d31 ;
    rdfs:seeAlso sh-gufo:concernsRelationshipType-subjects-shape ;
    .

kb:TemporaryRelationshipSituation-2a64c96d-9629-4c29-9d63-e6f69cdbdabd
    a gufo:TemporaryRelationshipSituation ;
    gufo:concernsRelatedEndurant kb:Endurant-7ac20c26-1f62-40c1-a6e0-0e96fb088e9b ;
    rdfs:seeAlso sh-gufo:concernsRelatedEndurant-subjects-shape ;
    .
\end{lstlisting}

\paragraph{Failure tests --- A6.2 and A6.3.}

\begin{lstlisting}[style=turtlestyle]
kb:TemporaryRelationshipSituation-7a1b5c6d-9e0f-4a1b-3c4d-5e6f7a8b9c0d
    a gufo:TemporaryRelationshipSituation ;
    gufo:concernsRelatedEndurant kb:Endurant-7ac20c26-1f62-40c1-a6e0-0e96fb088e9b ;
    rdfs:comment "This node is expected to trigger a validation error for
        having no gufo:concernsRelationshipType triple."@en ;
    rdfs:seeAlso sh-gufo:TemporaryRelationshipSituation-shape ;
    .

kb:TemporaryRelationshipSituation-8b2c6d7e-0f1a-4b2c-4d5e-6f7a8b9c0d1e
    a gufo:TemporaryRelationshipSituation ;
    gufo:concernsRelationshipType kb:RelationshipType-0813ed87-2f47-4a1b-a063-834b467c3d31 ;
    rdfs:comment "This node is expected to trigger a validation error for
        having no gufo:concernsRelatedEndurant triple."@en ;
    rdfs:seeAlso sh-gufo:TemporaryRelationshipSituation-shape ;
    .
\end{lstlisting}

\paragraph{pytest.}

\begin{lstlisting}[style=pythonstyle]
def test_exemplar_xfail_validation_temporary_relationship_situation() -> None:
    """
    Verifies A6.2 and A6.3: minCount 1 on both concernsRelationshipType and
    concernsRelatedEndurant. Two focus nodes expected, one per missing property.
    """
    validation_graph = Graph()
    validation_graph.parse("exemplars_XFAIL_validation.ttl")
    ns_kb = Namespace("http://example.org/kb/")
    ns_gufo = Namespace("http://purl.org/nemo/gufo#")
    pairs: Set[Tuple[URIRef, URIRef]] = set()
    for result in validation_graph.query("""\
PREFIX sh: <http://www.w3.org/ns/shacl#>
PREFIX sh-gufo: <http://example.org/shapes/sh-gufo/>
SELECT ?nSituation ?nPath
WHERE {
  ?nValidationResult
    a sh:ValidationResult ;
    sh:sourceShape sh-gufo:TemporaryRelationshipSituation-shape ;
    sh:sourceConstraintComponent sh:MinCountConstraintComponent ;
    sh:focusNode ?nSituation ;
    sh:resultPath ?nPath ;
    .
}
"""):
        assert isinstance(result, ResultRow)
        pairs.add((result[0], result[1]))
    assert pairs == {
        (ns_kb["TemporaryRelationshipSituation-7a1b5c6d-9e0f-4a1b-3c4d-5e6f7a8b9c0d"],
         ns_gufo["concernsRelationshipType"]),
        (ns_kb["TemporaryRelationshipSituation-8b2c6d7e-0f1a-4b2c-4d5e-6f7a8b9c0d1e"],
         ns_gufo["concernsRelatedEndurant"]),
    }
\end{lstlisting}

%% ───────────────────────────────────────────────────────────────────────────
\subsection{Rule B1 --- QualityValueAttributionSituation Uses Exactly One Quality Value Pattern}
\label{sec:rule-b1}
%% ───────────────────────────────────────────────────────────────────────────

\subsubsection*{(i) Rule Description}

A \gufoclass{QualityValueAttributionSituation} (abbreviated QVAS) must be
associated with a quality value using \emph{exactly one} of two mutually
exclusive patterns: either a reified quality value via
\gufoprop{concernsReifiedQualityValue}, or a literal value via
\gufoprop{concernsQualityValue}. The two patterns are alternatives---using both
simultaneously is semantically inconsistent, and using neither makes the
situation meaningless. The gUFO OWL implementation expresses this as
\nolinkurl{owl:unionOf} of two cardinality restrictions, one per pattern. The
correct SHACL operator is \shprop{xone}.

\textbf{Sources:} \texttt{formatted-gufo.ttl}; gUFO documentation \S2.5.

\subsubsection*{(ii) Test Aspects}

\begin{tabular}{lL{9.5cm}}
  \toprule
  \textbf{ID} & \textbf{Aspect} \\
  \midrule
  B1.1 & Completeness: a QVAS with neither property is invalid \\
  B1.2 & Mutual exclusion: a QVAS with both properties simultaneously is invalid \\
  B1.3 & Upper bound on \gufoprop{concernsReifiedQualityValue}: at most one \\
  B1.4 & Upper bound on \gufoprop{concernsQualityValue}: at most one \\
  B1.5 & Range of \gufoprop{concernsReifiedQualityValue}: must be a \gufoclass{QualityValue} \\
  B1.6 & Domain of \gufoprop{concernsReifiedQualityValue}: must be a QVAS \\
  B1.7 & Domain of \gufoprop{concernsQualityValue}: must be a QVAS \\
  B1.8 & Range of \gufoprop{concernsQualityValue}: must be a \texttt{Literal} \\
  \bottomrule
\end{tabular}

\subsubsection*{(iii) DL Expressions}

\textbf{B1.1 + B1.2 --- Exclusive disjunction (\texttt{xone} semantics):}
\begin{multline*}
  \mathsf{QVAS} \sqsubseteq
  \bigoplus\bigl\{
    (\exists\,\mathsf{concernsReifiedQualityValue}.\mathsf{QualityValue}),\\
    \hspace{6em}(\exists\,\mathsf{concernsQualityValue}.\top)
  \bigr\}
\end{multline*}
where $\bigoplus$ denotes exclusive disjunction (exactly one disjunct holds).

\textbf{B1.3:}
\begin{equation*}
  \mathsf{QVAS} \sqsubseteq {\leq}1\;\mathsf{concernsReifiedQualityValue}.\mathsf{QualityValue}
\end{equation*}

\textbf{B1.4:}
\begin{equation*}
  \mathsf{QVAS} \sqsubseteq {\leq}1\;\mathsf{concernsQualityValue}.\top
\end{equation*}

\textbf{B1.5:}
$\top \sqsubseteq \forall\,\mathsf{concernsReifiedQualityValue}.\mathsf{QualityValue}$

\textbf{B1.6:}
$\exists\,\mathsf{concernsReifiedQualityValue}.\top \sqsubseteq \mathsf{QVAS}$

\textbf{B1.7:}
$\exists\,\mathsf{concernsQualityValue}.\top \sqsubseteq \mathsf{QVAS}$

\textbf{B1.8:}
$\top \sqsubseteq \forall\,\mathsf{concernsQualityValue}.\mathsf{Literal}$

\subsubsection*{(iv) Implementation}

\paragraph{SHACL shapes.}

\begin{lstlisting}[style=turtlestyle]
sh-gufo:QualityValueAttributionSituation-shape
    a sh:NodeShape ;
    sh:xone (
        [
            a sh:NodeShape ;
            sh:property [
                a sh:PropertyShape ;
                sh:path gufo:concernsReifiedQualityValue ;
                sh:qualifiedMaxCount "1"^^xsd:integer ;
                sh:qualifiedMinCount "1"^^xsd:integer ;
                sh:qualifiedValueShape [ a sh:NodeShape ; sh:class gufo:QualityValue ; ] ;
            ] ;
        ]
        [
            a sh:NodeShape ;
            sh:property [
                a sh:PropertyShape ;
                sh:maxCount "1"^^xsd:integer ;
                sh:minCount 1 ;
                sh:path gufo:concernsQualityValue ;
            ] ;
        ]
    ) ;
    sh:property
        [
            a sh:PropertyShape ;
            sh:path gufo:concernsQualityType ;
            sh:qualifiedMaxCount "1"^^xsd:integer ;
            sh:qualifiedValueShape [ a sh:NodeShape ; sh:class gufo:EndurantType ; ] ;
        ] ,
        [
            a sh:PropertyShape ;
            sh:path [ sh:inversePath gufo:standsInQualifiedAttribution ] ;
            sh:qualifiedMaxCount "1"^^xsd:integer ;
            sh:qualifiedValueShape [ a sh:NodeShape ; sh:class gufo:Endurant ; ] ;
        ] ;
    sh:targetClass gufo:QualityValueAttributionSituation ;
    .
\end{lstlisting}

\paragraph{Success tests.}

\begin{lstlisting}[style=turtlestyle]
kb:QualityValueAttributionSituation-20fec4d2-6e44-49ed-94aa-799845277a8c
    a gufo:QualityValueAttributionSituation ;
    gufo:concernsQualityValue "" ;
    rdfs:seeAlso
        sh-gufo:QualityValueAttributionSituation-shape ,
        sh-gufo:concernsQualityValue-subjects-shape ;
    .

kb:QualityValueAttributionSituation-87d3fdff-f584-43f9-9339-974d46cc3187
    a gufo:QualityValueAttributionSituation ;
    gufo:concernsReifiedQualityValue kb:QualityValue-bbbc153a-8979-4ca8-a1cc-89a53194c5c8 ;
    rdfs:seeAlso sh-gufo:concernsReifiedQualityValue-subjects-shape ;
    .
\end{lstlisting}

\paragraph{Failure tests --- B1.1 and B1.2.}

\begin{lstlisting}[style=turtlestyle]
# B1.2: both patterns simultaneously
kb:QualityValueAttributionSituation-9c3d7e8f-1a2b-4c3d-5e6f-7a8b9c0d1e2f
    a gufo:QualityValueAttributionSituation ;
    gufo:concernsReifiedQualityValue kb:QualityValue-bbbc153a-8979-4ca8-a1cc-89a53194c5c8 ;
    gufo:concernsQualityValue "" ;
    rdfs:comment "This node is expected to trigger a validation error for
        having both gufo:concernsReifiedQualityValue and gufo:concernsQualityValue
        simultaneously."@en ;
    rdfs:seeAlso sh-gufo:QualityValueAttributionSituation-shape ;
    .

# B1.1: neither pattern present
kb:QualityValueAttributionSituation-0d4e8f9a-2b3c-4d4e-6f7a-8b9c0d1e2f3a
    a gufo:QualityValueAttributionSituation ;
    rdfs:comment "This node is expected to trigger a validation error for
        having neither gufo:concernsReifiedQualityValue nor
        gufo:concernsQualityValue."@en ;
    rdfs:seeAlso sh-gufo:QualityValueAttributionSituation-shape ;
    .
\end{lstlisting}

\paragraph{pytest.}

\begin{lstlisting}[style=pythonstyle]
def test_exemplar_xfail_validation_quality_value_attribution_xone() -> None:
    """
    Verifies B1.1 and B1.2: sh:xone on QualityValueAttributionSituation-shape.
    Both the all-absent and the both-present cases must produce an
    XoneConstraintComponent violation.
    """
    validation_graph = Graph()
    validation_graph.parse("exemplars_XFAIL_validation.ttl")
    ns_kb = Namespace("http://example.org/kb/")
    n_focus_nodes: Set[URIRef] = set()
    for result in validation_graph.query("""\
PREFIX sh: <http://www.w3.org/ns/shacl#>
PREFIX sh-gufo: <http://example.org/shapes/sh-gufo/>
SELECT ?nSituation
WHERE {
  ?nValidationResult
    a sh:ValidationResult ;
    sh:sourceShape sh-gufo:QualityValueAttributionSituation-shape ;
    sh:sourceConstraintComponent sh:XoneConstraintComponent ;
    sh:focusNode ?nSituation ;
    .
}
"""):
        assert isinstance(result, ResultRow)
        assert isinstance(result[0], URIRef)
        n_focus_nodes.add(result[0])
    assert n_focus_nodes == {
        ns_kb["QualityValueAttributionSituation-9c3d7e8f-1a2b-4c3d-5e6f-7a8b9c0d1e2f"],
        ns_kb["QualityValueAttributionSituation-0d4e8f9a-2b3c-4d4e-6f7a-8b9c0d1e2f3a"],
    }
\end{lstlisting}

%% ───────────────────────────────────────────────────────────────────────────
\subsection{Rule B2 --- concernsQualityType Range Must Be a Subclass of gufo:Quality}
\label{sec:rule-b2}
%% ───────────────────────────────────────────────────────────────────────────

\subsubsection*{(i) Rule Description}

The property \gufoprop{concernsQualityType} identifies the quality type whose
value is attributed in a QVAS. gUFO specifies that the range must be a subclass
of \gufoclass{Quality}---not merely any \gufoclass{EndurantType}. Without this
tighter constraint, a \gufoclass{Kind} subclassing \gufoclass{Object} would
incorrectly pass validation as a quality type.

\textbf{Sources:} gUFO documentation for \texttt{concernsQualityType};
\texttt{formatted-gufo.ttl}.

\subsubsection*{(ii) Test Aspects}

\begin{tabular}{lL{9.5cm}}
  \toprule
  \textbf{ID} & \textbf{Aspect} \\
  \midrule
  B2.1 & Range is an \gufoclass{EndurantType} \\
  B2.2 & Range is specifically a subclass of \gufoclass{Quality} \\
  B2.3 & Domain: subject must be a QVAS \\
  B2.4 & Upper bound: at most one \gufoprop{concernsQualityType} per QVAS \\
  \bottomrule
\end{tabular}

\subsubsection*{(iii) DL Expressions}

\textbf{B2.1:}
$\top \sqsubseteq \forall\,\mathsf{concernsQualityType}.\mathsf{EndurantType}$

\textbf{B2.2:}
\begin{equation*}
  \top \sqsubseteq \forall\,\mathsf{concernsQualityType}.
  \bigl(\mathsf{EndurantType} \sqcap \exists\,\mathsf{rdfs{:}subClassOf}^{+}.\mathsf{Quality}\bigr)
\end{equation*}

\textbf{B2.3:}
$\exists\,\mathsf{concernsQualityType}.\top \sqsubseteq \mathsf{QVAS}$

\textbf{B2.4:}
$\mathsf{QVAS} \sqsubseteq {\leq}1\;\mathsf{concernsQualityType}.\mathsf{EndurantType}$

\subsubsection*{(iv) Implementation}

\paragraph{SHACL shapes.}

\begin{lstlisting}[style=turtlestyle]
sh-gufo:concernsQualityType-objects-shape
    a sh:NodeShape ;
    sh:message "The object of gufo:concernsQualityType must be an
        EndurantType that subclasses gufo:Quality."@en ;
    sh:class gufo:EndurantType ;
    sh:nodeKind sh:BlankNodeOrIRI ;
    sh:property [
        a sh:PropertyShape ;
        sh:hasValue gufo:Quality ;
        sh:message "The object of gufo:concernsQualityType must be a
            subclass of gufo:Quality."@en ;
        sh:path [ sh:oneOrMorePath rdfs:subClassOf ] ;
    ] ;
    sh:targetObjectsOf gufo:concernsQualityType ;
    .

sh-gufo:concernsQualityType-subjects-shape
    a sh:NodeShape ;
    sh:class gufo:QualityValueAttributionSituation ;
    sh:property [
        a sh:PropertyShape ;
        sh:maxCount "1"^^xsd:integer ;
        sh:path gufo:concernsQualityType ;
    ] ;
    sh:targetSubjectsOf gufo:concernsQualityType ;
    .
\end{lstlisting}

\paragraph{Success test.}

\begin{lstlisting}[style=turtlestyle]
kb:QualityType-3a0ff32d-9743-46e4-bdcd-e86777f98551
    a gufo:Kind , owl:Class ;
    rdfs:subClassOf gufo:Quality ;
    rdfs:seeAlso sh-gufo:concernsQualityType-objects-shape ;
    .
\end{lstlisting}

\paragraph{Failure test --- B2.2.}

\begin{lstlisting}[style=turtlestyle]
kb:Kind-1e2f9a0b-3c4d-4e1f-7a8b-9c0d1e2f3a4b
    a gufo:Kind , owl:Class ;
    rdfs:subClassOf gufo:Object ;
    rdfs:comment "This node is expected to trigger a validation error for
        being used as gufo:concernsQualityType while not subclassing
        gufo:Quality."@en ;
    rdfs:seeAlso sh-gufo:concernsQualityType-objects-shape ;
    .

kb:QualityValueAttributionSituation-1e2f0a1b-4d5e-4f2a-8b9c-0d1e2f3a4b5c
    a gufo:QualityValueAttributionSituation ;
    gufo:concernsQualityValue "" ;
    gufo:concernsQualityType kb:Kind-1e2f9a0b-3c4d-4e1f-7a8b-9c0d1e2f3a4b ;
    rdfs:comment "This node uses a non-Quality EndurantType as
        gufo:concernsQualityType."@en ;
    rdfs:seeAlso sh-gufo:concernsQualityType-objects-shape ;
    .
\end{lstlisting}

\paragraph{pytest.}

\begin{lstlisting}[style=pythonstyle]
def test_exemplar_xfail_validation_concerns_quality_type_range() -> None:
    """Verifies B2.2: object of concernsQualityType must subclass gufo:Quality."""
    validation_graph = Graph()
    validation_graph.parse("exemplars_XFAIL_validation.ttl")
    ns_kb = Namespace("http://example.org/kb/")
    n_focus_nodes: Set[URIRef] = set()
    for result in validation_graph.query("""\
PREFIX sh: <http://www.w3.org/ns/shacl#>
PREFIX sh-gufo: <http://example.org/shapes/sh-gufo/>
SELECT ?nType
WHERE {
  ?nValidationResult
    a sh:ValidationResult ;
    sh:sourceShape sh-gufo:concernsQualityType-objects-shape ;
    sh:sourceConstraintComponent sh:HasValueConstraintComponent ;
    sh:focusNode ?nType ;
    .
}
"""):
        assert isinstance(result, ResultRow)
        assert isinstance(result[0], URIRef)
        n_focus_nodes.add(result[0])
    assert n_focus_nodes == {
        ns_kb["Kind-1e2f9a0b-3c4d-4e1f-7a8b-9c0d1e2f3a4b"],
    }
\end{lstlisting}

%% ───────────────────────────────────────────────────────────────────────────
\subsection{Rule B3 --- MaterialRelationshipType Must Derive From a Subclass of gufo:Relator}
\label{sec:rule-b3}
%% ───────────────────────────────────────────────────────────────────────────

\subsubsection*{(i) Rule Description}

A \gufoclass{MaterialRelationshipType} represents a relationship type whose
instances are derived from a relator. According to \citet{fonseca2019}, a
material relationship type must be derived (via \gufoprop{isDerivedFrom}) from a
subclass of \gufoclass{Relator}---not merely any subclass of
\gufoclass{ExtrinsicAspect}. Using an \gufoclass{ExtrinsicMode} subclass as the
derivation basis is the \emph{RelRig} anti-pattern \citep{sales2015}, which
incorrectly treats a mode-mediated relation as a relator-mediated one.

\textbf{Sources:} \citet{fonseca2019}; \citet{sales2015}; \texttt{formatted-gufo.ttl}.

\subsubsection*{(ii) Test Aspects}

\begin{tabular}{lL{9.5cm}}
  \toprule
  \textbf{ID} & \textbf{Aspect} \\
  \midrule
  B3.1 & Domain: subjects of \gufoprop{isDerivedFrom} must be a
         \gufoclass{MaterialRelationshipType} or \gufoclass{ComparativeRelationshipType} \\
  B3.2 & Objects of \gufoprop{isDerivedFrom} must be an \gufoclass{EndurantType}
         subclassing \gufoclass{Aspect} \\
  B3.3 & \gufoclass{ComparativeRelationshipType} derives from a subclass of
         \gufoclass{IntrinsicAspect} \\
  B3.4 & \gufoclass{MaterialRelationshipType} derives from a subclass of
         \gufoclass{Relator} specifically \\
  \bottomrule
\end{tabular}

\subsubsection*{(iii) DL Expressions}

\textbf{B3.1 --- Domain:}
\begin{equation*}
  \exists\,\mathsf{isDerivedFrom}.\top
  \sqsubseteq \mathsf{MaterialRelationshipType} \sqcup \mathsf{ComparativeRelationshipType}
\end{equation*}

\textbf{B3.2 --- Object subclasses Aspect:}
\begin{equation*}
  \top \sqsubseteq \forall\,\mathsf{isDerivedFrom}.
  \bigl(\mathsf{EndurantType} \sqcap \exists\,\mathsf{rdfs{:}subClassOf}^{+}.\mathsf{Aspect}\bigr)
\end{equation*}

\textbf{B3.3 --- ComparativeRelationshipType derives from IntrinsicAspect:}
\begin{multline*}
  \mathsf{ComparativeRelationshipType} \sqsubseteq \forall\,\mathsf{isDerivedFrom}.\\
  \bigl(\mathsf{EndurantType} \sqcap \exists\,\mathsf{rdfs{:}subClassOf}^{+}.\mathsf{IntrinsicAspect}\bigr)
\end{multline*}

\textbf{B3.4 --- MaterialRelationshipType derives from Relator:}
\begin{multline*}
  \mathsf{MaterialRelationshipType} \sqsubseteq \forall\,\mathsf{isDerivedFrom}.\\
  \bigl(\mathsf{EndurantType} \sqcap \exists\,\mathsf{rdfs{:}subClassOf}^{+}.\mathsf{Relator}\bigr)
\end{multline*}

\subsubsection*{(iv) Implementation}

\paragraph{SHACL shapes.}

\begin{lstlisting}[style=turtlestyle]
sh-gufo:isDerivedFrom-MaterialRelationshipType-subjects-shape
    a sh:NodeShape ;
    rdfs:seeAlso <https://nemo-ufes.github.io/gufo/#isDerivedFrom> ;
    sh:message "gufo:MaterialRelationshipTypes should be derived from a subclass
        of gufo:Relator ( https://nemo-ufes.github.io/gufo/#isDerivedFrom )."@en ;
    sh:or (
        [
            a sh:NodeShape ;
            sh:not [
                a sh:NodeShape ;
                sh:class gufo:MaterialRelationshipType ;
            ] ;
        ]
        [
            a sh:NodeShape ;
            sh:property [
                a sh:PropertyShape ;
                sh:path gufo:isDerivedFrom ;
                sh:property [
                    a sh:PropertyShape ;
                    sh:hasValue gufo:Relator ;
                    sh:path [ sh:oneOrMorePath rdfs:subClassOf ] ;
                ] ;
            ] ;
        ]
    ) ;
    sh:targetSubjectsOf gufo:isDerivedFrom ;
    .
\end{lstlisting}

\paragraph{Success tests.}

\begin{lstlisting}[style=turtlestyle]
kb:MaterialRelationshipType-2133b940-8a06-4856-bf46-684592243af4
    a owl:ObjectProperty , gufo:MaterialRelationshipType ;
    gufo:isDerivedFrom kb:EndurantType-959b6cff-b5a4-4168-be1f-d66d0c8033d6 ;
    rdfs:seeAlso
        sh-gufo:isDerivedFrom-MaterialRelationshipType-subjects-shape ,
        sh-gufo:isDerivedFrom-subjects-shape ;
    .

kb:EndurantType-959b6cff-b5a4-4168-be1f-d66d0c8033d6
    a gufo:Kind , owl:Class ;
    rdfs:subClassOf gufo:Relator ;  # subclasses Relator, not just ExtrinsicAspect
    rdfs:seeAlso
        sh-gufo:isDerivedFrom-MaterialRelationshipType-subjects-shape ,
        sh-gufo:isDerivedFrom-objects-shape ;
    .
\end{lstlisting}

\paragraph{Failure test --- B3.4.}

\begin{lstlisting}[style=turtlestyle]
kb:EndurantType-2f3a0b1c-5e6f-4a2b-8c9d-0e1f2a3b4c5d
    a gufo:Kind , owl:Class ;
    rdfs:subClassOf gufo:ExtrinsicMode ;
    rdfs:comment "This node subclasses gufo:ExtrinsicMode (not gufo:Relator)
        and triggers an invalid isDerivedFrom for a MaterialRelationshipType."@en ;
    rdfs:seeAlso sh-gufo:isDerivedFrom-MaterialRelationshipType-subjects-shape ;
    .

kb:MaterialRelationshipType-3a4b1c2d-6f7a-4b3c-9d0e-1f2a3b4c5d6e
    a owl:ObjectProperty , gufo:MaterialRelationshipType ;
    gufo:isDerivedFrom kb:EndurantType-2f3a0b1c-5e6f-4a2b-8c9d-0e1f2a3b4c5d ;
    rdfs:comment "This node is expected to trigger a validation error for
        deriving from a subclass of gufo:ExtrinsicMode rather than
        gufo:Relator."@en ;
    rdfs:seeAlso sh-gufo:isDerivedFrom-MaterialRelationshipType-subjects-shape ;
    .
\end{lstlisting}

\paragraph{pytest.}

\begin{lstlisting}[style=pythonstyle]
def test_exemplar_xfail_validation_material_relationship_type_relator() -> None:
    """Verifies B3.4: MaterialRelationshipType must derive from a Relator subclass."""
    validation_graph = Graph()
    validation_graph.parse("exemplars_XFAIL_validation.ttl")
    ns_kb = Namespace("http://example.org/kb/")
    n_focus_nodes: Set[URIRef] = set()
    for result in validation_graph.query("""\
PREFIX sh: <http://www.w3.org/ns/shacl#>
PREFIX sh-gufo: <http://example.org/shapes/sh-gufo/>
SELECT ?nRelType
WHERE {
  ?nValidationResult
    a sh:ValidationResult ;
    sh:sourceShape sh-gufo:isDerivedFrom-MaterialRelationshipType-subjects-shape ;
    sh:sourceConstraintComponent sh:OrConstraintComponent ;
    sh:focusNode ?nRelType ;
    .
}
"""):
        assert isinstance(result, ResultRow)
        assert isinstance(result[0], URIRef)
        n_focus_nodes.add(result[0])
    assert n_focus_nodes == {
        ns_kb["MaterialRelationshipType-3a4b1c2d-6f7a-4b3c-9d0e-1f2a3b4c5d6e"],
    }
\end{lstlisting}

%% ───────────────────────────────────────────────────────────────────────────
\subsection{Rule C1 --- EndurantType Satisfies Two Independent Exclusivity Constraints}
\label{sec:rule-c1}
%% ───────────────────────────────────────────────────────────────────────────

\subsubsection*{(i) Rule Description}

Every \gufoclass{EndurantType} must satisfy two \emph{independent} exclusive
disjunctions simultaneously: it must be exactly one of \gufoclass{NonRigidType}
or \gufoclass{RigidType} (the rigidity partition), and exactly one of
\gufoclass{NonSortal} or \gufoclass{Sortal} (the sortality partition). In SHACL
these are expressed as two separate \shprop{xone} constraints on the same shape
node. Because SHACL conjoins multiple property values on a node shape, both
constraints must be satisfied independently---this is \emph{not} a single
disjunction but a conjunction of two exclusive disjunctions.

\textbf{Sources:} \citet{guizzardi2018}; \texttt{formatted-gufo.ttl}
\nolinkurl{owl:disjointUnionOf} on \texttt{EndurantType}.

\subsubsection*{(ii) Test Aspects}

\begin{tabular}{lL{9.5cm}}
  \toprule
  \textbf{ID} & \textbf{Aspect} \\
  \midrule
  C1.1 & Every \gufoclass{EndurantType} is exactly one of
         \gufoclass{NonRigidType} or \gufoclass{RigidType} \\
  C1.2 & Every \gufoclass{EndurantType} is exactly one of
         \gufoclass{NonSortal} or \gufoclass{Sortal} \\
  \bottomrule
\end{tabular}

\subsubsection*{(iii) DL Expressions}

\textbf{C1.1 --- Rigidity partition:}
\begin{equation*}
  \mathsf{EndurantType} \sqsubseteq
  \bigoplus\{\mathsf{NonRigidType},\;\mathsf{RigidType}\}
\end{equation*}

\textbf{C1.2 --- Sortality partition:}
\begin{equation*}
  \mathsf{EndurantType} \sqsubseteq
  \bigoplus\{\mathsf{NonSortal},\;\mathsf{Sortal}\}
\end{equation*}

\textbf{Combined --- both must hold:}
\begin{equation*}
  \mathsf{EndurantType} \sqsubseteq
  \bigoplus\{\mathsf{NonRigidType},\;\mathsf{RigidType}\}
  \;\sqcap\;
  \bigoplus\{\mathsf{NonSortal},\;\mathsf{Sortal}\}
\end{equation*}

\subsubsection*{(iv) Implementation}

\paragraph{SHACL shape.}

\begin{lstlisting}[style=turtlestyle]
sh-gufo:EndurantType-shape
    a sh:NodeShape ;
    rdfs:comment "This shape applies two independent sh:xone constraints.
        The first requires every EndurantType to be exactly one of
        (NonRigidType, RigidType). The second requires every EndurantType to
        be exactly one of (NonSortal, Sortal). Both constraints must hold
        simultaneously and independently -- this is a conjunction of two
        exclusive disjunctions, not a single disjunction."@en ;
    sh:targetClass gufo:EndurantType ;
    sh:xone
        (
            [ a sh:NodeShape ; sh:class gufo:NonRigidType ; ]
            [ a sh:NodeShape ; sh:class gufo:RigidType ; ]
        ) ,
        (
            [ a sh:NodeShape ; sh:class gufo:NonSortal ; ]
            [ a sh:NodeShape ; sh:class gufo:Sortal ; ]
        )
        ;
    .
\end{lstlisting}

\paragraph{Success test.}
Any \gufoclass{EndurantType} exemplar satisfies this shape; for example,
\gufoclass{Category} is a \gufoclass{RigidType} and a \gufoclass{NonSortal},
satisfying both \shprop{xone} constraints.

\begin{lstlisting}[style=turtlestyle]
kb:EndurantType-1e1f2313-7cbb-47ae-bdd0-bc4a4a5a654e
    a gufo:Category , owl:Class ;
    gufo:hasAssociatedQualityValueType
        kb:AbstractIndividualType-73fdd434-b356-4b08-8ad6-e112a8727e61 ;
    rdfs:seeAlso sh-gufo:hasAssociatedQualityValueType-subjects-shape ;
    .
\end{lstlisting}

\paragraph{Failure test.}

\begin{lstlisting}[style=turtlestyle]
kb:EndurantType-f6898a59-a3ec-4ed9-9b1d-0cd9a19d2663
    a gufo:EndurantType , owl:Class ;
    rdfs:comment "This is expected to trigger a validation error for not
        satisfying the disjoint-union, exactly-one clauses in
        EndurantType review."@en ;
    rdfs:seeAlso sh-gufo:EndurantType-shape ;
    .
\end{lstlisting}

\paragraph{pytest.}

\begin{lstlisting}[style=pythonstyle]
def test_exemplar_xfail_validation_enduranttype() -> None:
    """
    Verifies C1.1 and C1.2: bare EndurantType without rigidity/sortality subtype
    must trigger XoneConstraintComponent on EndurantType-shape.
    """
    validation_graph = Graph()
    validation_graph.parse("exemplars_XFAIL_validation.ttl")
    ns_kb = Namespace("http://example.org/kb/")
    n_focus_nodes: Set[URIRef] = set()
    for result in validation_graph.query("""\
PREFIX sh: <http://www.w3.org/ns/shacl#>
PREFIX sh-gufo: <http://example.org/shapes/sh-gufo/>
SELECT ?nEndurantType
WHERE {
  ?nValidationResult
    a sh:ValidationResult ;
    sh:sourceShape sh-gufo:EndurantType-shape ;
    sh:sourceConstraintComponent sh:XoneConstraintComponent ;
    sh:focusNode ?nEndurantType ;
    .
}
"""):
        assert isinstance(result, ResultRow)
        assert isinstance(result[0], URIRef)
        n_focus_nodes.add(result[0])
    assert n_focus_nodes == {
        ns_kb["EndurantType-f6898a59-a3ec-4ed9-9b1d-0cd9a19d2663"],
    }
\end{lstlisting}

%% ───────────────────────────────────────────────────────────────────────────
\subsection{Rule D1 --- QualityValueAttributionSituation Must Be Used by Exactly One Endurant}
\label{sec:rule-d1}
%% ───────────────────────────────────────────────────────────────────────────

\subsubsection*{(i) Rule Description}

A QVAS must be linked to exactly one \gufoclass{Endurant} via the inverse of
\gufoprop{standsInQualifiedAttribution}. The gUFO OWL implementation declares
\nolinkurl{owl:qualifiedCardinality "1"} on this inverse path, expressing that the
situation must be referenced by exactly one endurant---neither zero nor more
than one.

\textbf{Sources:} \texttt{formatted-gufo.ttl} \nolinkurl{owl:qualifiedCardinality "1"}
on inverse of \texttt{standsInQualifiedAttribution}.

\subsubsection*{(ii) Test Aspects}

\begin{tabular}{lL{9.5cm}}
  \toprule
  \textbf{ID} & \textbf{Aspect} \\
  \midrule
  D1.1 & Upper bound: at most one \gufoclass{Endurant} uses this situation via
         \gufoprop{standsInQualifiedAttribution} \\
  D1.2 & Lower bound: at least one \gufoclass{Endurant} must use this situation \\
  D1.3 & Domain of \gufoprop{standsInQualifiedAttribution}: subject must be
         an \gufoclass{Endurant} \\
  D1.4 & Range of \gufoprop{standsInQualifiedAttribution}: object must be a QVAS \\
  \bottomrule
\end{tabular}

\subsubsection*{(iii) DL Expressions}

\textbf{D1.1:}
$\mathsf{QVAS} \sqsubseteq {\leq}1\;\mathsf{standsInQualifiedAttribution}^{-}.\mathsf{Endurant}$

\textbf{D1.2:}
$\mathsf{QVAS} \sqsubseteq {\geq}1\;\mathsf{standsInQualifiedAttribution}^{-}.\top$

\textbf{D1.3:}
$\exists\,\mathsf{standsInQualifiedAttribution}.\top \sqsubseteq \mathsf{Endurant}$

\textbf{D1.4:}
$\top \sqsubseteq \forall\,\mathsf{standsInQualifiedAttribution}.\mathsf{QVAS}$

\subsubsection*{(iv) Implementation}

\paragraph{SHACL shapes.}

\begin{lstlisting}[style=turtlestyle]
sh-gufo:QualityValueAttributionSituation-shape
    a sh:NodeShape ;
    ...
    sh:property
        ...
        [
            a sh:PropertyShape ;
            sh:message "Every gufo:QualityValueAttributionSituation must be
                used by at least one Endurant via
                gufo:standsInQualifiedAttribution."@en ;
            sh:path [ sh:inversePath gufo:standsInQualifiedAttribution ] ;
            sh:qualifiedMaxCount "1"^^xsd:integer ;
            sh:qualifiedMinCount "1"^^xsd:integer ;
            sh:qualifiedValueShape [ a sh:NodeShape ; sh:class gufo:Endurant ; ] ;
        ] ;
    sh:targetClass gufo:QualityValueAttributionSituation ;
    .

sh-gufo:standsInQualifiedAttribution-subjects-shape
    a sh:NodeShape ;
    sh:class gufo:Endurant ;
    sh:targetSubjectsOf gufo:standsInQualifiedAttribution ;
    .

sh-gufo:standsInQualifiedAttribution-objects-shape
    a sh:NodeShape ;
    sh:class gufo:QualityValueAttributionSituation ;
    sh:nodeKind sh:BlankNodeOrIRI ;
    sh:targetObjectsOf gufo:standsInQualifiedAttribution ;
    .
\end{lstlisting}

\paragraph{Success tests.}

\begin{lstlisting}[style=turtlestyle]
kb:Endurant-30842dd2-146e-4720-81eb-ff8f655058a5
    a gufo:Object ;
    gufo:standsInQualifiedAttribution
        kb:QualityValueAttributionSituation-71781915-ac36-4264-83f4-3a89df7de700 ;
    rdfs:seeAlso sh-gufo:standsInQualifiedAttribution-subjects-shape ;
    .

kb:QualityValueAttributionSituation-71781915-ac36-4264-83f4-3a89df7de700
    a gufo:QualityValueAttributionSituation ;
    rdfs:seeAlso sh-gufo:standsInQualifiedAttribution-objects-shape ;
    .
\end{lstlisting}

\paragraph{Failure test --- D1.2.}

\begin{lstlisting}[style=turtlestyle]
kb:QualityValueAttributionSituation-2f3a1b2c-5d6e-4f3a-9b0c-1d2e3f4a5b6c
    a gufo:QualityValueAttributionSituation ;
    gufo:concernsQualityValue "" ;
    rdfs:comment "This node is expected to trigger a validation error for
        being a QVAS not referenced by any Endurant via
        gufo:standsInQualifiedAttribution."@en ;
    rdfs:seeAlso sh-gufo:QualityValueAttributionSituation-shape ;
    .
\end{lstlisting}

\paragraph{pytest.}

\begin{lstlisting}[style=pythonstyle]
def test_exemplar_xfail_validation_quality_value_attribution_stands_in() -> None:
    """Verifies D1.2: qualifiedMinCount 1 on inverse standsInQualifiedAttribution."""
    validation_graph = Graph()
    validation_graph.parse("exemplars_XFAIL_validation.ttl")
    ns_kb = Namespace("http://example.org/kb/")
    n_focus_nodes: Set[URIRef] = set()
    for result in validation_graph.query("""\
PREFIX sh: <http://www.w3.org/ns/shacl#>
PREFIX sh-gufo: <http://example.org/shapes/sh-gufo/>
SELECT ?nSituation
WHERE {
  ?nValidationResult
    a sh:ValidationResult ;
    sh:sourceShape sh-gufo:QualityValueAttributionSituation-shape ;
    sh:sourceConstraintComponent sh:QualifiedMinCountConstraintComponent ;
    sh:focusNode ?nSituation ;
    .
}
"""):
        assert isinstance(result, ResultRow)
        assert isinstance(result[0], URIRef)
        n_focus_nodes.add(result[0])
    assert n_focus_nodes == {
        ns_kb["QualityValueAttributionSituation-2f3a1b2c-5d6e-4f3a-9b0c-1d2e3f4a5b6c"],
    }
\end{lstlisting}

%% ───────────────────────────────────────────────────────────────────────────
\subsection{Rule D2 --- TemporaryRelationshipSituation Must Be Used by Exactly One Endurant}
\label{sec:rule-d2}
%% ───────────────────────────────────────────────────────────────────────────

\subsubsection*{(i) Rule Description}

A \gufoclass{TemporaryRelationshipSituation} must be linked to exactly one
\gufoclass{Endurant} via the inverse of \gufoprop{standsInQualifiedRelationship}.
The gUFO OWL implementation declares \nolinkurl{owl:qualifiedCardinality "1"} on
this inverse path.

\textbf{Sources:} \nolinkurl{formatted-gufo.ttl} \nolinkurl{owl:qualifiedCardinality "1"} on inverse of \texttt{standsInQualifiedRelationship}.

\subsubsection*{(ii) Test Aspects}

\begin{tabular}{lL{9.5cm}}
  \toprule
  \textbf{ID} & \textbf{Aspect} \\
  \midrule
  D2.1 & Upper bound: at most one \gufoclass{Endurant} uses this situation via
         \gufoprop{standsInQualifiedRelationship} \\
  D2.2 & Lower bound: at least one \gufoclass{Endurant} must use this situation \\
  D2.3 & Domain of \gufoprop{standsInQualifiedRelationship}: subject must be
         an \gufoclass{Endurant} \\
  D2.4 & Range of \gufoprop{standsInQualifiedRelationship}: object must be a
         \gufoclass{TemporaryRelationshipSituation} \\
  \bottomrule
\end{tabular}

\subsubsection*{(iii) DL Expressions}

\textbf{D2.1:}
$\mathsf{TRS} \sqsubseteq {\leq}1\;\mathsf{standsInQualifiedRelationship}^{-}.\mathsf{Endurant}$

\textbf{D2.2:}
$\mathsf{TRS} \sqsubseteq {\geq}1\;\mathsf{standsInQualifiedRelationship}^{-}.\top$

\textbf{D2.3:}
$\exists\,\mathsf{standsInQualifiedRelationship}.\top \sqsubseteq \mathsf{Endurant}$

\textbf{D2.4:}
$\top \sqsubseteq \forall\,\mathsf{standsInQualifiedRelationship}.\mathsf{TRS}$

\subsubsection*{(iv) Implementation}

\paragraph{SHACL shapes.}

\begin{lstlisting}[style=turtlestyle]
sh-gufo:TemporaryRelationshipSituation-shape
    a sh:NodeShape ;
    sh:property [
        a sh:PropertyShape ;
        sh:message "Every gufo:TemporaryRelationshipSituation must be
            referenced by at least one Endurant via
            gufo:standsInQualifiedRelationship."@en ;
        sh:path [ sh:inversePath gufo:standsInQualifiedRelationship ] ;
        sh:qualifiedMaxCount "1"^^xsd:integer ;
        sh:qualifiedMinCount "1"^^xsd:integer ;
        sh:qualifiedValueShape [ a sh:NodeShape ; sh:class gufo:Endurant ; ] ;
    ] ;
    sh:targetClass gufo:TemporaryRelationshipSituation ;
    .

sh-gufo:standsInQualifiedRelationship-subjects-shape
    a sh:NodeShape ;
    sh:class gufo:Endurant ;
    sh:targetSubjectsOf gufo:standsInQualifiedRelationship ;
    .

sh-gufo:standsInQualifiedRelationship-objects-shape
    a sh:NodeShape ;
    sh:class gufo:TemporaryRelationshipSituation ;
    sh:nodeKind sh:BlankNodeOrIRI ;
    sh:targetObjectsOf gufo:standsInQualifiedRelationship ;
    .
\end{lstlisting}

\paragraph{Success tests.}

\begin{lstlisting}[style=turtlestyle]
kb:Endurant-c38ea234-5d65-40d2-9ce6-9ca9d070a3e6
    a gufo:Object ;
    gufo:standsInQualifiedRelationship
        kb:TemporaryRelationshipSituation-29301400-6b45-4176-b570-f52c28cd2789 ;
    rdfs:seeAlso sh-gufo:standsInQualifiedRelationship-subjects-shape ;
    .

kb:TemporaryRelationshipSituation-29301400-6b45-4176-b570-f52c28cd2789
    a gufo:TemporaryRelationshipSituation ;
    rdfs:seeAlso sh-gufo:standsInQualifiedRelationship-objects-shape ;
    .
\end{lstlisting}

\paragraph{Failure test --- D2.2.}

\begin{lstlisting}[style=turtlestyle]
kb:TemporaryRelationshipSituation-4b5c2d3e-7f8a-4c4b-0d1e-2f3a4b5c6d7e
    a gufo:TemporaryRelationshipSituation ;
    gufo:concernsRelationshipType kb:RelationshipType-0813ed87-2f47-4a1b-a063-834b467c3d31 ;
    gufo:concernsRelatedEndurant kb:Endurant-7ac20c26-1f62-40c1-a6e0-0e96fb088e9b ;
    rdfs:comment "This node is expected to trigger a validation error for
        being a TRS not referenced by any Endurant via
        gufo:standsInQualifiedRelationship."@en ;
    rdfs:seeAlso sh-gufo:TemporaryRelationshipSituation-shape ;
    .
\end{lstlisting}

\paragraph{pytest.}

\begin{lstlisting}[style=pythonstyle]
def test_exemplar_xfail_validation_temporary_relationship_situation_stands_in() -> None:
    """Verifies D2.2: qualifiedMinCount 1 on inverse standsInQualifiedRelationship."""
    validation_graph = Graph()
    validation_graph.parse("exemplars_XFAIL_validation.ttl")
    ns_kb = Namespace("http://example.org/kb/")
    n_focus_nodes: Set[URIRef] = set()
    for result in validation_graph.query("""\
PREFIX sh: <http://www.w3.org/ns/shacl#>
PREFIX sh-gufo: <http://example.org/shapes/sh-gufo/>
SELECT ?nSituation
WHERE {
  ?nValidationResult
    a sh:ValidationResult ;
    sh:sourceShape sh-gufo:TemporaryRelationshipSituation-shape ;
    sh:sourceConstraintComponent sh:QualifiedMinCountConstraintComponent ;
    sh:focusNode ?nSituation ;
    .
}
"""):
        assert isinstance(result, ResultRow)
        assert isinstance(result[0], URIRef)
        n_focus_nodes.add(result[0])
    assert n_focus_nodes == {
        ns_kb["TemporaryRelationshipSituation-4b5c2d3e-7f8a-4c4b-0d1e-2f3a4b5c6d7e"],
    }
\end{lstlisting}

%% ═══════════════════════════════════════════════════════════════════════════
\section{Summary of Rules and Test Aspects}
%% ═══════════════════════════════════════════════════════════════════════════

\begin{longtable}{L{1.5cm}L{10cm}c}
  \toprule
  \textbf{Rule} & \textbf{Description} & \textbf{Aspects} \\
  \midrule
  \endhead
  A1 & Every Aspect inheres in exactly one ConcreteIndividual & 4 \\
  A2 & Every ExtrinsicMode has at least one externallyDependsOn & 3 \\
  A3 & TemporaryConstitutionSituation has exactly one concernsConstitutedEndurant & 5 \\
  A4 & TemporaryInstantiationSituation has exactly one concernsNonRigidType & 5 \\
  A5 & TemporaryParthoodSituation has exactly one concernsTemporaryWhole & 5 \\
  A6 & TemporaryRelationshipSituation has one RelationshipType and at least one Endurant & 9 \\
  B1 & QualityValueAttributionSituation uses exactly one quality value pattern & 8 \\
  B2 & concernsQualityType range is a subclass of gufo:Quality & 4 \\
  B3 & MaterialRelationshipType derives from a subclass of gufo:Relator & 4 \\
  C1 & EndurantType satisfies two independent rigidity and sortality partitions & 2 \\
  D1 & QualityValueAttributionSituation is used by exactly one Endurant & 4 \\
  D2 & TemporaryRelationshipSituation is used by exactly one Endurant & 4 \\
  \midrule
  \textbf{Total} & & \textbf{57} \\
  \bottomrule
\end{longtable}

%% ═══════════════════════════════════════════════════════════════════════════
\begin{thebibliography}{10}

\bibitem[Carvalho et al.(2017)]{carvalho2017}
Carvalho, V.~A., Almeida, J.~P.~A., Fonseca, C.~M., and Guizzardi, G. (2017).
\newblock Multi-level ontology-based conceptual modeling.
\newblock \emph{Data \& Knowledge Engineering}, 109:3--24.

\bibitem[Fonseca et al.(2019)]{fonseca2019}
Fonseca, C.~M., Porello, D., Guizzardi, G., Almeida, J.~P.~A., and
  Guarino, N. (2019).
\newblock Relations in ontology-driven conceptual modeling.
\newblock In \emph{Conceptual Modeling -- ER~2019}, Lecture Notes in Computer
  Science, vol.~11788, pp.~1--15. Springer.

\bibitem[Guizzardi(2005)]{guizzardi2005}
Guizzardi, G. (2005).
\newblock \emph{Ontological Foundations for Structural Conceptual Models}.
\newblock PhD thesis, University of Twente, Enschede, The Netherlands.

\bibitem[Guizzardi et al.(2018)]{guizzardi2018}
Guizzardi, G., Fonseca, C.~M., Benevides, A.~B., Almeida, J.~P.~A., Porello,
  D., and Sales, T.~P. (2018).
\newblock Endurant types in ontology-driven conceptual modeling: Towards
  OntoUML~2.0.
\newblock In \emph{Conceptual Modeling -- ER~2018}, Lecture Notes in Computer
  Science, vol.~11157, pp.~136--150. Springer.

\bibitem[Guizzardi et al.(2022)]{guizzardi2022}
Guizzardi, G., Almeida, J.~P.~A., Fonseca, C.~M., Sales, T.~P., and
  Porello, D. (2022).
\newblock {UFO}: Unified Foundational Ontology.
\newblock \emph{Applied Ontology}, 17(1):167--210.

\bibitem[pyshacl]{pyshacl}
RDFLib/pySHACL (2023).
\newblock \emph{pySHACL: A Python validator for SHACL}.
\newblock \url{https://github.com/RDFLib/pySHACL}.

\bibitem[Sales and Guizzardi(2015)]{sales2015}
Sales, T.~P. and Guizzardi, G. (2015).
\newblock Ontological anti-patterns: Empirically uncovered error-prone
  structures in ontology-driven conceptual models.
\newblock \emph{Data \& Knowledge Engineering}, 99:72--104.

\bibitem[Sales and Guizzardi(2019)]{sales2019}
Sales, T.~P. and Guizzardi, G. (2019).
\newblock Anti-patterns for taxonomic structures in ontology-driven conceptual
  models.
\newblock In \emph{CEUR Workshop Proceedings}, vol.~2519.

\end{thebibliography}

\end{document}
